%! Absolute Value Equations \& Inequalities — Unit 2, Lesson 2.2 — Warm-Up (Answer Key)
\documentclass[10pt]{article}
\usepackage{algebra2-article}
\usepackage{algebra2-key}   % pulls in -boxes; do NOT also load -boxes
\newcommand{\TallMath}[1]{$\displaystyle #1\rule[-1.4em]{0pt}{3.2em}$}
\newcommand{\numline}[2]{%
  \draw[->, semithick, charcoal] (#1-0.4,0)--(#2+0.4,0);
  \foreach \x in {#1,...,#2}{\draw[charcoal] (\x,0.10)--(\x,-0.10) node[below, font=\scriptsize]{$\x$};}}

\begin{document}

\pageheader{Unit 2, Lesson 2.2}{Warm-Up --- Answer Key}
\namedateperiod

\vspace{0.10in}

\begin{notesbox}{Getting Ready: Distance, Solving, and Number Lines}
\small Today an absolute value asks a \emph{distance} question, and you'll solve the linear equation
inside the bars \emph{twice}. These three skills are the tools you'll need.

\vspace{4pt}
\textbf{1. Absolute value is distance from zero.} Evaluate each.
\[
  |-7| = \textcolor{keyred}{\mathbf{7}} \qquad |4| = \textcolor{keyred}{\mathbf{4}} \qquad
  |{-3}| + |5| = \textcolor{keyred}{\mathbf{8}} \qquad -|{-6}| = \textcolor{keyred}{\mathbf{-6}}
\]
How many different numbers have an absolute value of $7$? \ans{two}\; Name them: \ans{$7$ and $-7$}

\vspace{6pt}
\textbf{2. Solve the two-step equation.} (You'll repeat this move once the bars split into two cases.)
\[
  2x - 3 = 7 \qquad\Rightarrow\qquad x = \textcolor{keyred}{\mathbf{5}}
  \hspace{1.2cm}
  2x - 3 = -7 \qquad\Rightarrow\qquad x = \textcolor{keyred}{\mathbf{-2}}
\]

\vspace{6pt}
\textbf{3. Graph on a number line, then name it in interval notation.} Shade $x \ge -1$.
\begin{center}
\begin{tikzpicture}[scale=0.55]
  \numline{-5}{5}
  \draw[very thick, forest, ->] (-1,0)--(5.4,0);
  \fill[forest] (-1,0) circle (3.5pt);
\end{tikzpicture}
\end{center}
Interval notation: \ans{$[-1,\infty)$} \qquad {\footnotesize (Use $[\,$ or $]$ for ``or equal to,'' $($ or $)$ otherwise.)}
\end{notesbox}

\begin{teachernote}
Item~1 plants the two-case idea: because $|x|$ is a \emph{distance}, two numbers ($7$ and $-7$) share
an absolute value --- that is exactly why $|ax+b|=c$ splits into two equations. Item~2 is the linear
solving from Lesson~2.1; today students do it once for the ``$+c$'' case and once for the ``$-c$''
case. In item~3, tie the \emph{closed} dot at $-1$ to the bracket in $[-1,\infty)$; $\infty$ always
takes a parenthesis.
\end{teachernote}

\end{document}
